% Fig 5:fig16 -- Estructura del prompt ABSA para sentimientos electorales.
\begin{tikzpicture}[
    node distance=0.18cm,
    every node/.style={font=\scriptsize, align=left, inner sep=5pt},
    block/.style={draw, rounded corners=2pt, minimum width=12cm, minimum height=0.8cm},
    role/.style={block, fill=blue!10},
    ctx/.style={block, fill=cyan!12},
    instr/.style={block, fill=yellow!15},
    fmt/.style={block, fill=orange!12},
    ex/.style={block, fill=green!12}
]
    \node[role]  (rol) {\textbf{1. ROL.} ``Eres un analista politico experto en discurso electoral peruano.''};
    \node[ctx,   below=of rol]  (ctx) {\textbf{2. CONTEXTO.} Eleccion presidencial Peru 2026, 40 candidatos en liza con sus aliases.};
    \node[instr, below=of ctx]  (ins) {\textbf{3. INSTRUCCION.} Extrae candidatos mencionados, sentimiento ([-1,1]), tema y aspectos.};
    \node[fmt,   below=of ins]  (fmt) {\textbf{4. FORMATO.} JSON estructurado validado con esquema Pydantic v2.};
    \node[ex,    below=of fmt]  (ex)  {\textbf{5. EJEMPLOS.} 4 ejemplos few-shot calibrados al espa\~nol politico peruano.};
\end{tikzpicture}
